\section{Общие сведения}\label{ACME.Intro}

Технология ACME предназначена для автоматизации процессов управления 
сертификатами, прежде всего их выпуска. Сторонами ACME являются клиент и 
сервер. Сервер представляет УЦ, переадресуя ему запросы клиентов. Клиент
представляет сторону или службу, которая использует сертификаты в своей 
работе, например, почтовый сервер или TLS-сервер.

Взаимодействие между клиентом и сервером выполняется с помощью веб-сервисов.
Работа с веб-сервисом состоит в обращении с запросами на поддерживаемые им 
коммуникационные узлы. Ответы используются либо непосредственно, либо участвуют 
в формировании запросов на другие узлы. В запросы и ответы ACME включаются 
объекты JSON~\cite{RFC8259}.

Основное направление использования ACME~--- выпуск сертификатов TLS-серверов. В
сертификате фиксируется факт контроля над одним или несколькими доменными
именами со стороны сервера~--- владельца сертификата. При выпуске
сертификата требуется доказать факт контроля, и ACME позволяет создавать такого
рода доказательства. Доказательства ACME могут касаться не только
доменного имени, но и других объектов, называемых идентификаторами.

В самом начале клиент регистрируется на сервере, открывая здесь свою учетную
запись. С учетной записью связывается пара ключей ЭЦП: личный и открытый. Клиент
использует личный ключ для подписи запросов на веб-узлы, сервер использует
соответствующий открытый ключ при проверке подписей.

Зарегистрированный клиент получает сертификат по следующей схеме.

\begin{enumerate}
\item
Формирует заявку на выпуск сертификата. Заявка содержит перечень
идентификаторов, контроль над которыми клиент будет доказывать, а также данные,
не попадающие в стандартный запрос на выпуск сертификата.

\item
Доказывает контроль над заявленными идентификаторами.

\item
Финализирует заявку, отправляя запрос на выпуск сертификата.

\item
Ожидает завершения обработки запроса и получает выпущенный сертификат.
\end{enumerate}

Для отзыва сертификата клиент обращается с подписанным запросом, который
содержит ссылку на целевой сертификат.
